% !TeX root = ../教材题库_全部条目_转写工作稿.tex
% Source: XJ-B1-C03-EX-P082-Y0618
% 第3章 / 3.2 函数的基本性质 / 教材 p.82 / PDF 第89页
% OCR未校对：公式与文字均可能有误；未完成数学TeX转写；禁止直接发给学生
% 原题及相邻原解截图（核对用）：../题目原图/XJ-B1-C03-EX-P082-Y0618-01.webp
% 本文件只在转写完成后作为题干使用；OCR 条目尚不是可靠题干。
\begin{bankitem}{3.2 例题 例题3}{第3章 / 3.2 函数的基本性质 / 教材 p.82 / PDF 第89页}{XJ-B1-C03-EX-P082-Y0618}{OCR未校对：公式与文字均可能有误；未完成数学TeX转写；禁止直接发给学生}
\begin{ocrdraft}
[教材p.82 / PDF第89页]
\par
0二0]若函数丆0）一十2（a一D工十2在区间（一
\par
。。，4）上单调递减，求实数a
\par
的取值范围．
\par
解因为二次函数丆0）一十2一1十2的图象的对称轴为直线一1一。
\par
且开口向上，所以函数丆）在区间（一。。，1一的上单调递减．
\par
又已知该函数在区间（一。。，4）上单调递减，则
\par
你能借助函数图象来
\par
1一a冫4，
\par
说明实数a的取值范围吗？
\par
即磊一3．
\par
故实数的取值范围为（一。。
\end{ocrdraft}
\end{bankitem}
